\hypertarget{chapter-28-c20-coroutines}{%
\section{CHAPTER 28: C20 COROUTINES}\label{chapter-28-c20-coroutines}}

\hypertarget{c20-coroutines}{%
\section{C++20 COROUTINES}\label{c20-coroutines}}

\hypertarget{asynchronous-programming}{%
\subsection{1. Asynchronous
Programming}\label{asynchronous-programming}}

Coroutines are functions that can suspend execution to be resumed later.
They enable writing async code that looks synchronous.

\hypertarget{keywords}{%
\subsubsection{1.1 Keywords}\label{keywords}}

\begin{itemize}
\tightlist
\item
  \passthrough{\lstinline!co\_await!}: Suspend execution until an
  awaited operation completes.
\item
  \passthrough{\lstinline!co\_yield!}: Suspend execution and return a
  value (generator).
\item
  \passthrough{\lstinline!co\_return!}: Complete execution and return a
  final value.
\end{itemize}

A function containing any of these is a coroutine.

\hypertarget{generators-the-python-style}{%
\subsection{2. Generators (The Python
Style)}\label{generators-the-python-style}}

(Note: \passthrough{\lstinline!std::generator!} arrived in C++23, but
the machinery exists in C++20).

\begin{lstlisting}[language={C++}]
#include <coroutine>
#include <iostream>

struct Generator {
    struct promise_type;
    using handle_type = std::coroutine_handle<promise_type>;
    handle_type h;

    Generator(handle_type h) : h(h) {}
    ~Generator() { if (h) h.destroy(); }

    struct promise_type {
        int current_value;
        Generator get_return_object() { return Generator{handle_type::from_promise(*this)}; }
        std::suspend_always initial_suspend() { return {}; }
        std::suspend_always final_suspend() noexcept { return {}; }
        std::suspend_always yield_value(int value) {
            current_value = value; return {};
        }
        void return_void() {}
        void unhandled_exception() { std::terminate(); }
    };

    bool next() { h.resume(); return !h.done(); }
    int value() const { return h.promise().current_value; }
};

Generator sequence(int start, int end) {
    for (int i = start; i < end; ++i) {
        co_yield i; // Suspend here
    }
}

int main() {
    auto gen = sequence(0, 5);
    while (gen.next()) {
        std::cout << gen.value() << " ";
    }
}
\end{lstlisting}

\hypertarget{tasks-co_await}{%
\subsection{\texorpdfstring{3. Tasks
(\texttt{co\_await})}{3. Tasks (co\_await)}}\label{tasks-co_await}}

Used for async I/O. (Requires a library like
\passthrough{\lstinline!cppcoro!} or custom implementation in C++20).

\begin{lstlisting}[language={C++}]
Task<int> fetch_data() {
    auto conn = co_await connect("db.local");
    auto result = co_await conn.query("SELECT * FROM users");
    co_return result.size();
}
\end{lstlisting}

\hypertarget{under-the-hood}{%
\subsection{4. Under the Hood}\label{under-the-hood}}

The compiler generates a state machine for the coroutine, allocating the
``coroutine frame'' (local variables, promise object, instruction
pointer) on the heap (usually).
